\begin{abstract}
We introduce \emph{Magic Gems}, a novel geometric representation of magic squares as three-dimensional polyhedra. 
By mapping an $n \times n$ magic square onto a centered coordinate grid and interpreting cell values as vertical displacements, 
we construct a point cloud in $\R^3$ whose convex hull defines the Magic Gem.
This construction reveals a fundamental connection between the algebraic constraints defining magic squares 
and the statistical structure of their geometric representation: 
we prove that an arrangement is a magic square if and only if the covariance between spatial coordinates 
and value assignments vanishes identically.

This zero-covariance characterization provides a new lens through which to view magic squares, 
connecting classical combinatorics to concepts from multivariate statistics and rigid body mechanics.
We demonstrate that the Magic Gem representation is invariant under the dihedral symmetry group $D_4$, 
yielding a canonical geometric object for each equivalence class of magic squares.
Furthermore, we establish that magic configurations correspond to critical points in the covariance landscape, 
with perturbation analysis confirming their isolated nature.

Computational experiments on orders $n = 3$, $4$, and $5$ validate our theoretical results 
and illuminate the structure of the arrangement space.
The framework suggests new approaches to classification problems for magic squares 
and offers geometric insight into their remarkable scarcity among all possible arrangements.

\medskip
\noindent\textbf{Keywords:} Magic squares, convex polyhedra, covariance, combinatorial geometry, moment of inertia
\end{abstract}
